\section{Token Balancing and Response Length}
\label{sec:normalization}

\subsection{How does loss averaging change domain weights?}

The MOPD objective in Section~\ref{sec:setting} averages token losses
within each response. Open-MOPD analyzes a different implementation:
averaging all response tokens together gives instruction following
approximately $1\%$ of the token weight despite supplying $20\%$ of
the prompts \citep{gao2026openmopddiagnosingfixingcapability}.
Length imbalance therefore depends on the averaging rule; it is not
intrinsic to the response-normalized MOPD objective.

For response $y_i$, write its mean token loss as
$\ell_i=|y_i|^{-1}\sum_t\ell_{i,t}$, using any of the per-position
losses in Section~\ref{sec:setting}. Let $\mathcal B_d$ index the
responses from domain $d$ in the batch. The intended domain weights
are $\lambda_d\geq0$, with $\sum_d\lambda_d=1$.
Table~\ref{tab:normalization} compares three averaging rules using
only these quantities.

\begin{table}[!htbp]
\centering
\small
\renewcommand{\arraystretch}{1.15}
\begin{tabular*}{\linewidth}{@{\extracolsep{\fill}}ll@{}}
\toprule
Averaging rule & Batch loss \\
\midrule
All tokens together
  & $\bigl(\sum_i |y_i|\ell_i\bigr)\big/\bigl(\sum_i |y_i|\bigr)$ \\
Tokens within each domain
  & $\sum_d\lambda_d\,
    \bigl(\sum_{i\in\mathcal B_d}|y_i|\ell_i\bigr)
    \big/\bigl(\sum_{i\in\mathcal B_d}|y_i|\bigr)$ \\
Responses within each domain
  & $\sum_d\lambda_d\,|\mathcal B_d|^{-1}
    \sum_{i\in\mathcal B_d}\ell_i$ \\
\bottomrule
\end{tabular*}
\caption{Loss averaging on the same rollout batch. Each domain has at least
one nonempty response. The last rule recovers an ordinary response mean
when $\lambda_d$ equals the domain's share of responses.}
\label{tab:normalization}
\end{table}

With equal numbers of prompts, a domain producing responses ten times
as long receives ten times the weight under the first rule. The second
rule fixes each domain's total weight to $\lambda_d$, as in Open-MOPD's
token-share balancing, but longer responses still count more within
that domain. The third rule also gives responses equal weight inside
each domain. Thus domain balancing and response-length normalization
address different effects.

The same distinction holds for gradients. Writing the response gradient
as $g_i=\nabla_\theta\ell_i$, the token-weighted gradient within a domain
equals its mean response gradient plus the covariance between response
length and $g_i$, divided by mean response length.
Appendix~\ref{app:normalization} gives the exact fixed-batch identity
and its short proof. This is a weighted-average identity, not a new
optimizer. Equal domain weights do not imply equal gradient magnitudes
or compatible gradient directions.

\subsection{Experiment: compare averaging rules on identical batches and online}

\paragraph{Shared setting.}
All three studies use a Qwen3-1.7B-Base student on mathematics, code,
instruction following, and science. Each domain uses its own
frozen Qwen3-1.7B teacher trained with reinforcement learning (RL).
The reference run uses equal domain prompt quotas
and weights, $64$ responses per update, one update per fresh rollout
batch, and a $4{,}096$-token response cap. The initial budget is at most
$500$ updates per run, subject to a throughput pilot.
Appendix~\ref{app:setup} records the shared setup and evaluation protocol.
These studies are ongoing; preliminary training observations appear in
Section~\ref{sec:density}.

\paragraph{Fixed-batch comparison.}
At early, intermediate, and late checkpoints, apply all three rules to
the same responses and corrected teacher top-$64$ losses. Record domain
token shares, response lengths, gradient norms and directions before
clipping, and the length--gradient covariance. Check the decomposition
and evaluate each proposed step on a common held-out response bank.
Accumulate the correct sums and denominators across micro-batches and
workers: averaging unequal micro-batch means can change the intended
loss. An unequal-quota check first matches $\lambda_d$ to prompt shares,
then tests uniform weights separately, so changes in the domain prior
are distinguishable from length weighting.
\missing{Loss-reduction audit and gradient measurements.}

\paragraph{Online comparison.}
Train three branches differing only in the averaging rule. They share
initialization, teachers, prompt schedule, length limits, and optimizer
settings; each generates responses from its own evolving student.
The response-averaged top-$64$ branch is reused in
Sections~\ref{sec:subnetworks} and~\ref{sec:density}.
Report per-domain capability, macro-average, and worst-domain change
against both updates and generated tokens, with GPU time reported
separately. A supporting $1{,}024$ versus $4{,}096$ cap comparison
records truncation and answer-completion rates. Shortening responses
changes the available prefixes and may remove solution endings, so it
is a separate intervention from loss reweighting. Neither the algebra
nor the cap comparison guarantees a capability improvement.
\missing{Per-domain capability curves for the three averaging rules.}
